\section{Five-dimensional inducing blocks: one remaining Gaussian group}
\label{sec:rank-five-inducing}

\paragraph{Scope and check status.}
This is a new reduction, recorded on 21 November 2026.  It is a
necessary-condition theorem, not a construction of a dominant family.
The scalar-field and Chevalley--Weil calculations have an exact finite
certificate.  The parent independently read the complete proof and
reran the finite certificate, recording PASS for both on 21 November
2026.  The G34 agent independently read D-R5-01--03, the cited
Cohen rank-five proof, and a separate field/spectrum certificate;
these checks also passed.  The later D-R5-04 computation and
geometric conclusion also passed an independent whole-proof review,
an independent all-eighty-class local-form calculation, and fresh
exact group enumeration by the G34 agent.  The parent separately
reran its exact group and local-form certificates.
The published classification input is Cohen's rank-five theorem and
normalizer results, with the 1978 correction; the real-boundary argument
uses Lewis--Reiner's published equal-pair extraction theorem.

\subsection*{D-R5-01. The reflection closure and the real boundary}

Suppose a finite-cover source has an exact imaginary quadratic character
field $K$, actual Weil sixfold signature $(3,3)$, period rank nine,
and nonsplit actual polarization.  Suppose its complex coefficient
representation is induced from an irreducible representation $\tau$
of dimension five.  Write $J$ for the faithful image of $\tau$ and
$L=\mathbf Q(\chi_\tau)$.  We prove that
\[
                         K\in\{\mathbf Q(i),\mathbf Q(\sqrt{-3})\}.
                                                               \tag{R5.1}
\]
Moreover, if $K\ne\mathbf Q(\sqrt{-3})$, then necessarily
\[
 \begin{gathered}
 K=\mathbf Q(i),\qquad L=\mathbf Q(\zeta_{12}),\qquad
 R=G_{33},\qquad J=\mu_4G_{33},\\
 |J|=103680,
 \end{gathered}                                                  \tag{R5.2}
\]
where $R$ is the normal subgroup generated by all reflections of $J$.
The intermediate rational base has exactly twelve branch points,
ten with reflection inertia and two with inertia of order four and
moving rank three.  The latter two spectra are respectively
\[
                 (1,1,-1,-1,i),\qquad(1,1,-1,-1,-i).
                                                               \tag{R5.3}
\]
These assertions allow arbitrary original representation degree and
do not assume that the inducing representation has a $K$-model.

If $\tau$ is imprimitive, its prime dimension makes its blocks
one-dimensional.  Transitivity of induction makes the original
representation monomial, and D-MON excludes a dominant nonsplit
source.  Thus $\tau$ is primitive.
Pass to the intermediate connected $J$-cover $C_0\to X$.
Frobenius reciprocity identifies its complex multiplicity variation
with the original one; the fixed polarization determines the
conjugate filtration, so the full original period rank factors
through this intermediate-cover moduli, as in D-BI-01.
If $h=g(X)$, the Euler equation is
\[
                            D=16-10h.
\]
For $h=1$ there are at most six branch points and at most six
marked-base parameters.  Larger $h$ is impossible.  Thus rank nine
forces $h=0$ and
\[
 r\ge12,\qquad D=16,\qquad
                  \sum_j\bigl(\operatorname{rank}(1-g_j)-1\bigr)
                       =16-r\le4.                           \tag{R5.4}
\]
In particular there are at least eight reflection inertias.

The reflection normal closure $R$ is nontrivial and irreducible.
Indeed its fixed space is $J$-stable and must vanish; the
nonorthogonality components of its reflection lines give an
orthogonal direct-sum decomposition permuted by $J$, which
primitivity forbids unless there is only one component.
If $R$ is imprimitive, Cohen's normalizer theorem
\cite[Proposition~2.13 and its proof, pp.~391--392]{FIELD:Cohen}
puts $J$ in a monomial group: none of its three exceptional
dimensions is five.  This again contradicts the already excluded
monomial case.

We next exclude every real $R$, without needing a list of real
rank-five Coxeter types.  In real orthogonal coordinates for $R$,
every unitary normalizer is scalar times a real orthogonal matrix,
by \cite[Lemma~2.12 and its proof, p.~391]{FIELD:Cohen} and
uniqueness of the positive invariant form.
There must be at least two inertia matrices with nonreal scalar
phase.  Otherwise all inertia matrices are real, or their
product is a nonreal scalar times a real matrix; the first case
would make $\chi_\tau$ and its induced character real, and the
second cannot have product one.
Such an inertia has fixed dimension at most two: the eigenspaces
of the corresponding real orthogonal matrix for its two distinct
conjugate nonreal eigenvalues have equal dimension.
Its moving rank is therefore at least three.
Equation (R5.4) now forces exactly twelve branch points, with ten
real reflections and two nonreal scalar inertias of rank three.
Hurwitz moves put the latter two at the end.  Write them as
\[
                 g_1=\lambda A_1,
                 \qquad g_2=\lambda^{-1}A_2,
                                                               \tag{R5.5}
\]
where $A_i$ are real orthogonal and $\lambda$ is nonreal.
Their original phases have real product $\pm1$, since
$g_1g_2\in R$; absorb the sign into $A_2$ to obtain (R5.5).

Put $w=(g_1g_2)^{-1}\in R$.  Lewis--Reiner
\cite[Corollary~1.4, p.~1459; proof \S4.3, pp.~1480--1481]{BI:LR}
changes the length-ten reflection factorization of $w$ by Hurwitz
moves into
\[
              (r_1,r_1,\ldots,r_q,r_q,s_1,\ldots,s_\ell),
              \qquad q=(10-\ell)/2,                         \tag{R5.6}
\]
with a shortest final factorization of length $\ell$.
Carter's independent-root criterion in
\cite[Proposition~4.1, p.~1473]{BI:LR} gives $\ell\le5$;
parity gives $\ell\in\{0,2,4\}$.
The telescoping product identity and independence of the successive
transformed roots show
\[
             \ker(w-1)=\bigcap_j\ker(s_j-1),
             \qquad \dim\ker(w-1)=5-\ell.
\]
Let $H_0$ be generated by the $s_j,g_1,g_2$, and let
$E=\ker(g_1-1)\cap\ker(g_2-1)$, of dimension $h_0$.
The preceding fixed-space identity gives $E=\tau^{H_0}$.
Both $E$ and $\overline E$ lie in the real fixed space of $w$;
they are disjoint because $A_1$ has the distinct eigenvalues
$\lambda^{-1}$ and $\lambda$ on them.  Consequently
\[
                         2h_0\le5-\ell.
                                                               \tag{R5.7}
\]

Bubble off the $q$ equal reflection pairs on rational leaves of
the base, each attached by a node of trivial inertia.  The source
dual graph is connected because the full tuple generates $J$.
Its edge module is $\mathbf C[J]^q$, and its vertex module is
$\mathbf C[J/H_0]\oplus\bigoplus_j\mathbf C[J/\langle r_j\rangle]$.
Taking $\tau$-multiplicities in the connected graph's cellular
sequence gives
\[
 \dim\operatorname{Hom}_J(\tau,H_1(\Gamma,\mathbf C))
       =5q-4q-h_0=q-h_0
       \ge\frac{10-\ell}{2}
                      -\left\lfloor\frac{5-\ell}{2}\right\rfloor
       =3.                                                    \tag{R5.8}
\]
Smooth all nodes with positive integer orders.  The positive
diagonal edge-length pairing is nondegenerate on the graph cycle
space, so the resulting single Picard--Lefschetz logarithm $N$
has exactly this rank on the coefficient multiplicity.  It has
$N^2=0$ on a six-dimensional space; hence its rank is at most
three, and (R5.8) forces rank three.

This is maximal isotropy over the original field $K$ as well.
The intermediate Hurwitz component has dimension nine, and the
original period rank nine makes its classifying map dominant.
Take a closed point of the generic fibre of that classifying map
and normalize its closure.  After shrinking it gives a finite
etale cover of a dense open of the intermediate component carrying
the original $K$-variation.  Choose the node-smoothing arc generically
so its punctured part lies in that open; after a finite base change
it lifts to the finite cover.  The logarithm on the original
$K$-variation is thus a positive integer multiple of $N$ after
the chosen complex embedding.  It is $K$-linear, square zero,
and has $K$-rank three.  Skew adjointness for the actual Hermitian
polarization gives
$h(Nx,Ny)=-h(x,N^2y)=0$.
Its image is a three-dimensional isotropic $K$-space, proving
splitness and contradicting the hypothesis.
This is the graph and finite-multisection argument of D-BI-12
with the rank-five numbers written out; it does not descend an
arbitrary $L$-isotropic subspace by assertion.

Therefore $R$ is complex primitive of rank five.  Cohen's complete
rank-five result identifies it with
\[
                             R=G_{33}=W(K_5)
                                                               \tag{R5.9}
\]
\cite[\S5.10 and proof, pp.~430--432; Table~4.16, p.~412]{FIELD:Cohen}.
Its reflection model is over $F=\mathbf Q(\omega)$, all its
reflections have order two, its degrees are $4,6,10,12,18$, and
its order is $51840$.  The model and these data appear in
\cite[\S4.15, pp.~409--410; \S4.16, pp.~411--413]{FIELD:Cohen}.
Its exact character field is $F$: it is contained in $F$, and a
complex reflection group of this type has nonreal character by
\cite[\S1.4(vi) and proof, p.~382]{FIELD:Cohen}.
Its scalar center is $\mu_2$, by the degree-gcd formula
\cite[\S1.4(ii), p.~381]{FIELD:Cohen}.

\subsection*{D-R5-02. Scalar-coupled subgroups and the exact branch profile}

The corrected scalar-normalizer result gives
\[
                       R\subset J\subset\mathbf C^*R
\]
\cite[Proposition~5.14, p.~435]{FIELD:Cohen}
\cite[p.~613]{FIELD:CohenErr}.
There is no unexamined subgroup-extension assumption here.
If $g=\lambda w\in J$, with $w\in R$, then
$\lambda I=gw^{-1}\in J$ because $R\subset J$.
Consequently $J=CR$, where $C$ is its finite cyclic scalar subgroup.
In the fixed $F$-model of $R$, put
\[
                 J_F=J\cap\operatorname{GL}_5(F)
                     =(C\cap F^*)R.
                                                               \tag{R5.10}
\]
The quotient $J/J_F$ is cyclic.  Membership in $J_F$ is independent
of the decomposition $g=\lambda w$, since its ambiguity is a scalar
in $R$, namely an element of $\mu_2\subset F$.

Induction of characters gives $K\subset L$.  Since $R\subset J$,
one also has $F\subset L$.  If $J=J_F$, its matrices are over $F$
and hence $L=F$, which forces $K=F$.
Otherwise at least two branch entries are outside $J_F$: they
generate a nontrivial quotient and their product there is one.
For such an entry write $g=\lambda w$, $\lambda\notin F$.
If it fixed at least three dimensions, the minimal polynomial of
$\lambda^{-1}$ over $F$ would occur at least three times in the
degree-five characteristic polynomial of $w$.  Its degree would
be at most one, a contradiction.  Thus each such entry moves at
least three dimensions.
Equation (R5.4) now forces
\[
 r=12,\quad \text{ten reflection entries},\quad
 \text{two entries outside }J_F\text{, each of moving rank three}.
                                                               \tag{R5.11}
\]
All other entries are in $R$.  Moving the two exceptional entries
to the end, their product lies in $R$; hence either one, say
$g=\lambda w$, generates $J$ together with $R$.
It follows that
\[
                         J=\langle R,\lambda I\rangle,
                         \qquad L=F(\lambda).
                                                               \tag{R5.12}
\]
The equality of fields uses both directions: all matrices are
over $F(\lambda)$, while $F\subset L$ and the scalar trace
$\operatorname{Tr}(\lambda I)=5\lambda$ puts $\lambda$ in $L$.

Now $\lambda^{-1}$ has multiplicity exactly two on $w$.
The same repeated-minimal-polynomial argument gives
$[F(\lambda):F]=2$.
If $d$ is the order of $\lambda$, then $\varphi(d)\le4$ and
\[
                         d\in\{1,2,3,4,5,6,8,10,12\}.
\]
The values $1,2,3,6$ lie in $F$ and are excluded.
For $d=5,10$, the sole quadratic subfield of $\mathbf Q(\zeta_d)$
is $\mathbf Q(\sqrt5)$; for $d=8$ its quadratic subfields are
$\mathbf Q(i),\mathbf Q(\sqrt2),\mathbf Q(\sqrt{-2})$.
None is $F$, so these three values would give extension degree
four over $F$ and are excluded.  Thus
\[
                    d\in\{4,12\},\qquad
                    L=F(i)=\mathbf Q(\zeta_{12}).
                                                               \tag{R5.13}
\]
Its only imaginary quadratic subfields are $\mathbf Q(i)$ and
$F$, proving (R5.1).

We can sharpen (R5.13).  The nontrivial automorphism of
$F(\lambda)/F$ sends $\lambda$ to $-\lambda$ in both remaining
cases.  Thus $w$ has eigenvalues
\[
             \lambda^{-1},\lambda^{-1},
             -\lambda^{-1},-\lambda^{-1},\mu,
             \qquad\mu\in F.
\]
Every element of $R$ has determinant $\varepsilon\in\{1,-1\}$,
because its generating reflections have order two.
Taking the product of these five eigenvalues gives
$\mu=\varepsilon\lambda^4$.
Therefore
\[
             \operatorname{Spec}(g)
                    =(1,1,-1,-1,\varepsilon\lambda^5).
                                                               \tag{R5.14}
\]
If $d=12$, then $g^2$ is a reflection whose nonidentity eigenvalue
$\lambda^{10}$ has order six.  It belongs to $J$ and hence to its
reflection normal closure $R$, contradicting the order-two
reflection list of $G_{33}$.  Thus $d=4$.
Equations (R5.12)--(R5.14) give $J=\mu_4R$ and an outside inertia
spectrum $(1,1,-1,-1,\pm i)$ of order four.  Its square is a
reflection in $R$.
The ten reflection determinants have product one, so the determinants
of the two outside inertias are reciprocal.  They must have the
opposite signs in (R5.3).

Since $R\cap\mu_4=\mu_2$, one has $J/R\simeq C_2$ and
$|J|=2\cdot51840=103680$.
The quotient cover $C_0/R\to\mathbf P^1$ has exactly two branch
points, both with full order-two inertia, and hence genus zero.
The local ages in (R5.3) are $5/4$ and $7/4$; the ten reflections
have age $1/2$ each.  Total age is $8=5+3$, as is the total dual
age, so this necessary profile has the required Chevalley--Weil
multiplicities $(3,3)$.
This verifies the signatures of a possible profile; it does not
construct a product-one generating tuple.

\subsection*{D-R5-03. What the Gaussian case still requires}

For $K=\mathbf Q(i)$, the preceding conclusions freeze the remaining
five-dimensional gate to (R5.2)--(R5.3).
The natural $L$-factor of a generic $J$-cover in a relevant
twelve-branch component is an abelian twelvefold $A_\tau$, with
$H^1(A_\tau)$ of rank six over $L$.
If it actually arises by induction from an original $K$-sixfold
$B$, the rational representation identity is
\[
 \operatorname{Ind}_H^G\operatorname{Res}_{L/\mathbf Q}\tau
       \simeq
       \left(\operatorname{Res}_{K/\mathbf Q}\rho\right)^{\oplus2}.
                                                               \tag{R5.15}
\]
Indeed scalar extension of the $K$-model of $\rho$ to $L$ is
$\operatorname{Ind}\tau$, and restriction back to $\mathbf Q$
has multiplicity $[L:K]=2$.  Frobenius reciprocity on rational
Hodge structures therefore gives, in the original family,
\[
                              A_\tau\sim B^2.
                                                               \tag{R5.16}
\]
This is an isogeny of variations after finite choices of primitive
projectors, not just a comparison of complex dimensions.
Since the original rank-nine family dominates the nine-dimensional
intermediate Hurwitz component, this would force its generic
$A_\tau$, after finite base change, to have an endomorphism algebra
containing $M_2(\mathbf Q(i))$.

No theorem proving or disproving that generic extra endomorphism
condition for every remaining component is asserted here.
Likewise the actual original $K$-polarization has not been shown
split in this complex $G_{33}$ case.  A proof of generic
$\operatorname{End}^0(A_\tau)=L$, or a suitable maximal-rank
$K$-linear boundary logarithm, would exclude the remaining case;
an actual source construction would instead have to verify the
isogeny/field descent and the period rank.  Neither follows from
the branch-count equality or the Chevalley--Weil profile alone.

\subsection*{D-R5-04. Exact local classes, a norm determinant, and one excluded component}

\paragraph{Status.}
This additional proof and its two exact computations were completed
on 21 November 2026.  The G34 agent independently checked the full
proof, all eighty local forms, and a fresh exact enumeration of the
51840-element group and its forty conjugacy classes, recording PASS.
The parent independently reran the group and local-form certificates.
It constructs a natural $L$-cover, not an original $K$-sixfold source.

Use the coordinates $(x_1,x_2,x_3,x_4,t)$ for
$(e_5-e_6)^\perp\subset\mathbf C^6$, where the last two original
coordinates both equal $t$.  The positive Hermitian form is
\[
                         H_0=\operatorname{diag}(1,1,1,1,2).
\]
After removing irrelevant scalar multiples from Cohen's five roots
in \cite[\S4.15, p.~409]{FIELD:Cohen}, choose the norm-two roots
\[
 \begin{split}
 v_1&=(0,0,0,0,1),\\
 v_2&=\tfrac12(-1,1+2\omega,1,1,1),\\
 v_3&=(1,-1,0,0,0),\quad v_4=(0,1,-1,0,0),\quad
 v_5=(0,0,1,-1,0).
 \end{split}
\]
Their Gram matrix and determinant are
\[
 H=\begin{pmatrix}
 2&1&0&0&0\\
 1&2&\omega&\omega^2&0\\
 0&\omega^2&2&-1&0\\
 0&\omega&-1&2&-1\\
 0&0&0&-1&2
 \end{pmatrix},\qquad \det H=2.
                                                               \tag{R5.17}
\]
In this basis the generating reflections are
$s_j=I-e_j H_{j,*}$; their entries lie in $\mathbf Z[\omega]$.
The exact program \texttt{g33/enumerate\_g33.cpp} enumerates their
matrix closure by right multiplication, checks closure and order
$51840$, and partitions it by conjugation by the same generators.
All arithmetic is integral in $\mathbf Z[\omega]$.
It finds forty conjugacy classes and exactly forty-five reflections.
Characteristic polynomials are computed by exact Newton identities;
the second program independently checks them by six exact determinant
evaluations per class.

There are exactly two classes with characteristic polynomials
\[
                 (x-1)(x^2+1)^2,
                 \qquad (x+1)(x^2+1)^2,
                                                               \tag{R5.18}
\]
each of size $540$ and order four.  A representative of the first is
the explicit word
\[
               w=s_1s_2s_1s_3s_2s_1s_4s_2s_5s_4.
                                                               \tag{R5.19}
\]
The existence is also visible without enumeration: Cohen's root list
contains the $D_4$ roots in the first four coordinates
\cite[pp.~409--410 and the $D_4$ stabilizer argument, p.~413]{FIELD:Cohen}.
The even signed permutation
\[
        j(x_1,x_2,x_3,x_4,t)=(x_2,-x_1,x_4,-x_3,t)
\]
lies in that $W(D_4)$ subgroup and has the first characteristic
polynomial in (R5.18).  Uniqueness of the enumerated class therefore
identifies its class with that of (R5.19).

Since adjoining the central scalar $i$ does not change conjugation
on $R$, the two cosets $R$ and $iR$ give exactly eighty conjugacy
classes in $J=\mu_4R$.  The program \texttt{g33/local\_forms.py}
works in the exact field $L=\mathbf Q(\omega,i)$ and, for every
class, computes a column basis $C$ of $\operatorname{im}(1-g)$,
its actual Hermitian Gram matrix $C^*HC$, and
$b(g)=\det(1-g|_{\operatorname{im}(1-g)})$.
It checks all eighty classes and finds exactly two outside classes
of moving rank three, namely those in (R5.18) multiplied by $i$.
There are no outside reflections.

For the two forced outside types, use the representatives $ij$ and
$-ij$.  Moving bases for them are, respectively,
\[
 (1,i,0,0,0),\ (0,0,1,i,0),\ (0,0,0,0,1)
\]
and the same vectors with $i$ replaced by $-i$.
Both moving Gram matrices are $\operatorname{diag}(2,2,2)$.
Consequently the determinant classes and denominators are
\[
 \det H_{\mathrm{mov}}=8,\qquad
 b(ij)=4(1-i),\qquad b(-ij)=4(1+i).
                                                               \tag{R5.20}
\]
The coordinate-column bases selected by the exact program give
determinants $16$ and $8(2+\sqrt3)$ instead; these agree with $8$
modulo norms because
\[
 N_{L/L^+}(1+i)=2,\qquad
 N_{L/L^+}(1-i\omega)=2+\sqrt3,
 \qquad L^+=\mathbf Q(\sqrt3).
\]
These norm equalities are checked in the field arithmetic as well.

The proof of D-DET-01 works verbatim for a CM quadratic extension
$L/L^+$: every occurrence of $\mathbf Q$ in its Hermitian determinant
argument is replaced by $L^+$, while complex conjugation is the CM
involution.  No geometric or linear-algebra step uses $L^+=\mathbf Q$.
Every reflection moving line has determinant class $2$, and its
$b$-factor is $2$.  For any product-one tuple with the forced
ten-plus-two profile, (R5.20) therefore gives signed parabolic
discriminant
\[
 \left[
 \frac{2^{10}\,8^2}
      {2^{10}\,4(1-i)\,4(1+i)}
 \right]
                  =[2]=[1]
             \quad\text{in }(L^+)^*/N_{L/L^+}(L^*).
                                                               \tag{R5.21}
\]
Thus no obstruction to rational determinant descent is obtained:
the determinant already has a rational representative which is a
Gaussian norm.

This conclusion does not descend splitness to $K=\mathbf Q(i)$.
For a concrete warning about that inference, the Gaussian Hermitian
space $U^2\perp\langle1,-3\rangle$ has signature $(3,3)$ and
nonsplit signed discriminant $3$, since $3$ is not a Gaussian norm.
After adjoining $\sqrt3$ it is split, because the last plane has
isotropic vector $(\sqrt3,1)$.  Hence even a split Hermitian space
over $L$ can be obtained from a nonsplit Gaussian space.

There is nevertheless an explicit natural $L$-cover and an exact
component obstruction.  With the word $w$ in (R5.19), set $g=iw$.
The twelve-entry tuple
\[
             (s_1,s_1,\ldots,s_5,s_5,g,g^{-1})
                                                               \tag{R5.22}
\]
has product one and generates $J$: its reflection entries generate
$R$, and $gw^{-1}=iI$ supplies the scalar.  The exact program
verifies the word, product and scalar extraction.
It has the required local spectra and coefficient multiplicities
$(3,3)$; its connected $J$-cover has genus
\[
 1+\frac{103680}{2}
        \left(-2+10\left(1-\frac12\right)
                    +2\left(1-\frac14\right)\right)
                 =233281.
\]
This does not produce a Gaussian sixfold: its natural rational
factor is the $L$-abelian twelvefold $A_\tau$ of D-R5-03.

Finally bubble off the five doubled reflection pairs in (R5.22).
The central component has inertia group $\langle g\rangle$, whose
fixed space on $\tau$ has dimension two.  The connected source
graph therefore has coefficient multiplicity
\[
                      5\cdot5-5\cdot4-2=3.
                                                               \tag{R5.23}
\]
Its positive node smoothing gives a square-zero logarithm of rank
three.  If a rank-nine original Gaussian source dominated this
particular intermediate Hurwitz component, the finite-multisection
argument in D-R5-01 would lift that logarithm to the original
$K$-variation and make its actual Hermitian space split.
Thus the component containing (R5.22) cannot supply a dominant
nonsplit original Gaussian sixfold.  No transitivity assertion
identifying every possible Nielsen component with this one is
made; that, or an independent generic-endomorphism calculation,
remains the narrower residual gate.
